\documentclass[11pt,letterpaper]{article}
\usepackage[spanish]{babel}
\usepackage{fontspec}
\setmainfont{TeX Gyre Heros}
\usepackage[margin=2.1cm]{geometry}
\usepackage{amsmath}
\usepackage{siunitx}
\usepackage{booktabs}
\usepackage{tabularx}
\usepackage{enumitem}
\usepackage[most]{tcolorbox}
\usepackage{circuitikz}
\usepackage{xcolor}
\usepackage{hyperref}
\definecolor{azul}{HTML}{173B57}
\definecolor{morado}{HTML}{6A3D8F}
\definecolor{gris}{HTML}{F1F4F6}
\hypersetup{
  colorlinks=true,
  linkcolor=azul,
  urlcolor=azul,
  pdftitle={Prototipo - Guía Experiencia 01 - Ley de Ohm},
  pdfauthor={Diego García}
}
\setlength{\parindent}{0pt}
\setlength{\parskip}{5pt}
\setlist{itemsep=2pt,topsep=3pt}
\newtcolorbox{aviso}{
  colback=red!5,colframe=red!65!black,title=\textbf{PROTOTIPO - NO PUBLICAR}
}
\newtcolorbox{farmacia}{
  colback=morado!7,colframe=morado,title=\textbf{Puente químico-farmacéutico}
}
\begin{document}

{\color{azul}\Large\bfseries Experiencia 01 - Ley de Ohm}\\[2pt]
{\large Medición de voltaje en una resistencia}\\[8pt]
\textbf{Laboratorio de Física Aplicada a las Ciencias Farmacéuticas}

\begin{aviso}
Este documento existe para probar la estructura y la carga no publicada en
Canvas. Los modelos, rangos y límites de los instrumentos aún deben verificarse.
\end{aviso}

\section*{Pregunta experimental}

¿La caída de voltaje \(V\) en una resistencia fija es proporcional a la
corriente \(I\) que la atraviesa dentro del intervalo seguro del montaje?

\section*{Objetivos}

\begin{itemize}
  \item conectar correctamente voltímetro, amperímetro y ohmímetro;
  \item registrar pares de corriente y voltaje;
  \item representar \(V\) en función de \(I\);
  \item estimar la resistencia a partir de la pendiente;
  \item decidir si los datos son compatibles con \(V=RI\).
\end{itemize}

\section*{Fundamento mínimo}

Para un elemento óhmico a temperatura aproximadamente constante,
\[
  V=RI.
\]
Si se grafica \(V\) en el eje vertical e \(I\) en el horizontal, la pendiente
del ajuste lineal estima \(R\). Un intercepto distinto de cero o una curvatura
sistemática deben discutirse; no se eliminan datos solo para mejorar el ajuste.

\section*{Montaje}

\begin{center}
\begin{circuitikz}[american currents,american voltages]
  \draw (0,0) to[battery1,l={Fuente DC}] (0,3)
        to[ammeter,l={$A$}] (3,3)
        to[R,l={$R$}] (3,0) -- (0,0);
  \draw (3,3) -- (5,3)
        to[voltmeter,l={$V$}] (5,0) -- (3,0);
\end{circuitikz}
\end{center}

El amperímetro se conecta \textbf{en serie}; el voltímetro,
\textbf{en paralelo}. La resistencia se mide con ohmímetro únicamente cuando
está aislada y el circuito se encuentra desenergizado.

\section*{Instrumentos y materiales por grupo}

\begin{tabularx}{\textwidth}{@{}p{3.3cm}X@{}}
\toprule
\textbf{Elemento} & \textbf{Uso y dato que debe confirmarse}\\
\midrule
Multímetro 1 & Voltaje DC; modelo, rango, resolución y exactitud pendientes.\\
Multímetro 2 & Corriente DC; modelo, rango, resolución, entrada y fusible pendientes.\\
Función ohmímetro & Medición previa de \(R\) con circuito desenergizado.\\
Fuente DC regulada & Rango y límite de corriente pendientes de prueba.\\
Resistencia fija & Valor nominal, tolerancia y potencia máxima pendientes.\\
Montaje & Placa o tablero, cables compatibles e interruptor si corresponde.\\
\bottomrule
\end{tabularx}

\section*{Seguridad y cuidado}

\begin{enumerate}
  \item Mantener la fuente apagada mientras se arma o modifica el circuito.
  \item Solicitar revisión antes de energizar.
  \item Nunca conectar el amperímetro en paralelo con la fuente.
  \item Nunca usar el modo resistencia en un circuito energizado.
  \item No superar los límites definidos tras la prueba piloto.
  \item Apagar inmediatamente ante calentamiento, olor o lectura inestable.
\end{enumerate}

\section*{Procedimiento}

\begin{enumerate}
  \item Identificar el valor nominal, tolerancia y potencia de la resistencia.
  \item Con la fuente desconectada, medir la resistencia con el ohmímetro.
  \item Construir el circuito del esquema y comprobar rangos y terminales.
  \item Pedir la verificación del docente o ayudante.
  \item Encender desde el ajuste mínimo de la fuente.
  \item Registrar al menos seis pares \((I,V)\) dentro del intervalo aprobado.
  \item No reconectar ningún instrumento con el circuito energizado.
  \item Volver la fuente a cero, apagarla y desmontar.
\end{enumerate}

\section*{Registro de datos}

\begin{tabularx}{\textwidth}{@{}cXXXX@{}}
\toprule
\textbf{Punto} & \textbf{\(I\) (unidad)} & \textbf{\(V\) (unidad)}
& \textbf{\(V/I\) (\(\Omega\))} & \textbf{Observación}\\
\midrule
1 & & & &\\[5mm]
2 & & & &\\[5mm]
3 & & & &\\[5mm]
4 & & & &\\[5mm]
5 & & & &\\[5mm]
6 & & & &\\[5mm]
\bottomrule
\end{tabularx}

\section*{Análisis}

\begin{enumerate}
  \item Construir el gráfico \(V\) versus \(I\), con unidades y barras de
  incertidumbre si la resolución de los equipos permite estimarlas.
  \item Ajustar \(V=mI+b\) y reportar \(m\), \(b\) y una medida de calidad.
  \item Interpretar \(m\) como estimación de resistencia.
  \item Comparar \(m\) con la medición del ohmímetro y el valor nominal.
  \item Evaluar si los datos son compatibles con un comportamiento óhmico.
\end{enumerate}

\begin{farmacia}
Los conductímetros usados con soluciones electrolíticas también estudian la
respuesta eléctrica y el transporte de carga. Sin embargo, requieren considerar
la geometría de la celda, la temperatura y la calibración. Explique por qué una
resistencia metálica fija es un punto de partida útil, pero no un modelo completo
de una solución farmacéutica.
\end{farmacia}

\section*{Producto}

Completar el template del reporte con datos, gráfico, ajuste, comparación y una
conclusión sustentada. Este prototipo no constituye todavía una entrega oficial.

\vspace{3pt}
{\small\textbf{Fuentes de construcción.}
Syllabus oficial, AE 04; distribución semana a semana, semana 15;
\texttt{1Clases.pdf}, páginas PDF 169-172; \texttt{martins-2011.pdf}, páginas
PDF 249-251.}

\end{document}
